\chapter{Lagrangian perturbation theory and initial conditions}\label{ch:lpt}
\chapterstatus{written}{2026-11-14}

Lagrangian perturbation theory (LPT) describes structure formation by the displacement of fluid elements from their
initial positions. It is the standard way to set up particle initial conditions, the analytic backbone of the COLA
method (\cref{ch:nbody-pm}), and a forward model for reconstruction and field-level inference. This chapter gives the
LPT equations and their solutions to third order, the arbitrary-order recursion, the de-aliasing of the nonlinear
terms, the augmented and spherical-collapse variants, redshift-space reconstruction, and glass particle loads.

\begin{modulebox}{LPT}
\file{LPT/DisplacementFields.h} (namespace \nsp{FML::COSMOLOGY::LPT}): 1LPT, 2LPT and 3LPT potentials and displacement
fields, the scale-dependent displacement, ALPT and spherical collapse, \code{compute_LPT_growth_factors_LCDM},
\code{_GeffGLCDM}; \file{LPT/Dealiasing.h}; \file{LPT/ArbitraryOrder.h}; \file{LPT/Reconstruction.h} (RSD removal);
\file{LPT/GrowthHistory.h}.\par
Documentation: \file{LPT/README.md}, \file{ARBITRARY_ORDER_DERIVATION.md}, \file{ARBITRARY_ORDER_API.md},
\file{ARBITRARY_ORDER_PLAN.md}, \file{SCIENCE_RESULTS.md}. Tests: \file{tests/test_lpt.cpp}, \file{LPT/tests}.
\end{modulebox}
\begin{modulebox}{GlassIC}
\file{GlassIC/GlassIC.h}, \file{GlassIC/LoadIO.h} (namespace \nsp{FML::GLASSIC}). Documentation:
\file{GlassIC/README.md}, \file{GlassIC/tests/RESULTS.md}.
\end{modulebox}

\section{The Lagrangian equations of motion}\label{sec:lpt-equations}

A fluid element that starts at the Lagrangian position $\vq$ is at the Eulerian position
\begin{equation}
\vx(\vq,a)=\vq+\bm\Psi(\vq,a).
\end{equation}
With the deformation tensor $F_{ij}=\delta_{ij}+A_{ij}$, $A_{ij}=\partial\Psi_i/\partial q_j$, and its determinant
$J=\det F$, mass conservation is simply $1+\delta(\vx)=1/J(\vq)$. The Eulerian divergence of the equation of motion
$\ddot\vx+2H\dot\vx=-\nabla_x\Phi/a^2$ (comoving coordinates, $\nabla^2_x\Phi=\frac32\Om_mH^2a^2\delta$ in GR) becomes,
after multiplying by $J$ and with the cofactor matrix $C_{ij}=\partial J/\partial F_{ij}$,
\begin{equation}
\sum_{ij}C_{ij}\left(\partial_t^2+2H\partial_t\right)A_{ij}=\frac32\Om_m(a)H^2\,(J-1).
\label{eq:lpt-longitudinal}
\end{equation}
The initial flow is irrotational, which in Lagrangian coordinates is the Cauchy invariant: $F^T\,\dd F/\dd t$ is
symmetric. Expanding $\bm\Psi=\sum_n\bm\Psi^{(n)}$ order by order in the linear density gives, at each order, a
Poisson-type equation for the longitudinal part $\nabla_q\cdot\bm\Psi^{(n)}$ and, from third order, an equation for
the transverse part $\nabla_q\times\bm\Psi^{(n)}$, sourced by products of lower orders \cite{Bouchet1995,Matsubara2015}.

\section{First, second and third order}\label{sec:lpt-123}

Each order factorises into a growth factor (\cref{ch:growth-transfer}) and a spatial field, which \fml{} stores as
potentials: $\bm\Psi^{(1)}=\nabla\phi^{(1)}$, $\bm\Psi^{(2)}=\nabla\phi^{(2)}$,
$\bm\Psi^{(3)}=\nabla(\phi^{(3a)}+\phi^{(3b)})+\nabla\times\bm A^{(3)}$. With $H_{ij}=\partial_i\partial_j\phi^{(1)}$,
\begin{equation}
Q(A,B)=\Tr A\,\Tr B-\Tr(AB),\qquad K_i(A,B)=\epsilon_{ijk}\sum_lA_{lj}B_{lk},
\end{equation}
and $\nabla^2\chi=\frac12Q(H,H)=\sum_{i>j}(H_{ii}H_{jj}-H_{ij}^2)$, the potentials returned by
\code{compute_3LPT_potential_fourier} are, at the time of the input density $\delta$ (in units of the box),
\begin{align}
\nabla^2\phi^{(1)}&=-\delta, & \phi^{(2)}&=-\tfrac37\chi,\label{eq:lpt-potentials}\\
\nabla^2\phi^{(3a)}&=-\tfrac13\det H, & \nabla^2\phi^{(3b)}&=\tfrac{5}{21}Q\!\left(H,\nabla\nabla\chi\right),\qquad
\nabla^2\bm A^{(3)}=\tfrac17\bm K\!\left(H,\nabla\nabla\chi\right),\nonumber
\end{align}
which include the Einstein--de Sitter amplitudes of the growing modes (the Zel'dovich approximation is the first
order alone; $-3/7$ is the familiar 2LPT factor \cite{Scoccimarro1998,Crocce2006}). In EdS the $n$-th order grows as
$D^n$, so at another time and its velocity
\begin{equation}
\bm\Psi=D\bm\Psi^{(1)}+D^2\bm\Psi^{(2)}+D^3\bm\Psi^{(3)},\qquad
\frac{\dd\bm\Psi}{\dd t}=\frac{\dd\ln D}{\dd t}\left(D\bm\Psi^{(1)}+2D^2\bm\Psi^{(2)}+3D^3\bm\Psi^{(3)}\right),
\label{eq:lpt-velocity}
\end{equation}
which is what \code{compute_3LPT_displacement_field} returns (with $D=1$ at the time of $\delta$ and the units of
\code{dlogDdt} setting those of the velocity). In other cosmologies and theories of gravity the user, or
SimulationCode, multiplies each order by its own growth factor $D_1$, $D_2$, $D_{3a}$, $D_{3b}$ relative to the EdS
amplitude and uses the rates $f_i$; \cls{GrowthHistory} provides $D_1$, $D_2$ and their rates for this (and for the
$N$-body time stepping), with the peculiar velocity $v=a^2E f_1D_1\Psi^{(1)}+\ldots$ in the code velocity
$v=a^3E\,\dd x/\dd a$.

\paragraph*{Scale-dependent growth.} When the linear growth depends on scale (massive neutrinos, $f(R)$)
\code{from_LPT_potential_to_displacement_vector_scaledependent} applies a ratio $D(k,a)/D(k,a_{\rm ini})$ to the
potential in Fourier space before taking the gradient. Higher-order kernels with scale-dependent growth are
approximated by the EdS shape with the scale-dependent amplitude.

\begin{implementation}[Implementation: low-order LPT on the grid]
\begin{itemize}
\item All derivatives are spectral: $\partial_i\to ik_i$, $\nabla^{-2}\to-1/k^2$, the DC mode set to zero.
\item \emph{Nyquist convention}: for an even mesh the Nyquist mode $n_i=N_{\rm mesh}/2$ is its own mirror image, so
      odd ($ik_i$) and mixed ($k_ik_j$, $i\ne j$) derivatives of it depend on a sign choice and break reflection
      symmetry. They are set to zero (the derivative of the band-limited field there), while $k_i^2$ is kept; the
      results are then covariant under $x_i\to-x_i$, as in \file{TidalTensor/Hessian.h}. The same kernel is the
      continuous gradient of the PM force.
\item \emph{Memory}: the 3LPT construction keeps the first-order Hessian and streams the second-order Hessian in two
      batches through $N$ work grids: about 17 full grids at peak in 3D with all potentials and the curl, before the
      caller's outputs and FFT workspace.
\item 2D is supported (the vector potential then has only a $z$ component, $\nabla\times(A\hat z)=(\partial_yA,-\partial_xA)$).
\end{itemize}
\end{implementation}

\begin{cppcode}[Displacement fields to third order]
#include <FML/LPT/DisplacementFields.h>
using namespace FML::COSMOLOGY::LPT;

// delta_fourier: the linear density at the initial time (Fourier space, box units)
std::array<FML::GRID::FFTWGrid<3>, 3> Psi, dPsidt;
double dlogDdt = H_ini * f_ini;                          // sets the units of dPsidt
compute_3LPT_displacement_field<3>(delta_fourier, Psi, dPsidt, dlogDdt,
                                   /*ignore_curl=*/false, Dealiasing::Padding);
// Or the potentials, to scale each order with its own growth factor:
FML::GRID::FFTWGrid<3> phi1, phi2, phi3a, phi3b;  std::array<FML::GRID::FFTWGrid<3>, 3> A;
compute_3LPT_potential_fourier<3>(delta_fourier, phi1, phi2, phi3a, phi3b, A, false, Dealiasing::Padding);
\end{cppcode}

\section{De-aliasing the nonlinear terms}\label{sec:lpt-dealiasing}

The second- and third-order sources are products of fields evaluated in real space. On a grid the product of two
band-limited fields has modes up to twice (three times for cubic terms) the original bandwidth, which fold back onto
the resolved modes. With \code{Dealiasing::Padding} the normalised Fourier coefficients are embedded in a larger work
grid, the sources and potentials are computed there, and the results are cropped to the original mesh
\cite{Michaux2021}: at least $3/2$ of the width in every dimension for 2LPT (the quadratic rule) and at least twice
for 3LPT. The intermediate 2LPT potential keeps its full bandwidth on the work mesh while the third-order products are
formed; cropping it early would lose physical mode couplings. Work widths are rounded up to an even number divisible by
the number of MPI tasks; the spectral remapping exchanges only retained coefficients with chunked messages and never
replicates the grid. Only wave vectors with $2|n_i|<N_{\rm mesh}$ in every dimension are copied, so even-mesh
Nyquist planes are zeroed (also in the returned 1LPT potential); there is no radial cutoff or smoothing of the
sub-Nyquist modes. The cost is substantial: in 3D the work grids are $3.375\times$ (2LPT) and $8\times$ (3LPT) larger.
The default is \code{Dealiasing::None}, which keeps the original calculation. SimulationCode, ALPT and the
non-Gaussian generator do not enable it automatically.

\section{Arbitrary order}\label{sec:lpt-arbitrary}

\file{LPT/ArbitraryOrder.h} computes the growing-mode EdS solution to a user-selected order, with longitudinal and
transverse parts. With $g=a/a_{\rm ref}$ and no factorials,
$\bm\Psi(\vq,a)=\sum_{n\ge1}g^n\bm\Psi_n(\vq)$, $\nabla\cdot\bm\Psi_1=-\delta_{\rm ref}$. In EdS
$(\partial_t^2+2H\partial_t)g^n=H^2\lambda_ng^n$ with $\lambda_n=n(n+\frac12)$, so at order $n\ge2$ the unknown appears
in \cref{eq:lpt-longitudinal} only as $(\lambda_n-\frac32)\nabla\cdot\bm\Psi_n$. Matsubara's recursion
\cite{Matsubara2015} (Sec.~IV, Eqs.~62--64, EdS limit) gives, with $A^{(r)}=\nabla\bm\Psi_r$ and
$R(A,B,C)=\epsilon_{ijk}\epsilon_{pqr}A_{ip}B_{jq}C_{kr}$,
\begin{align}
d_n\equiv\nabla\cdot\bm\Psi_n&=-\frac12\sum_{\substack{r+s=n\\ r\le s}}M_{rs}
\left[1-\frac{4rs}{(2n+3)(n-1)}\right]Q(A^{(r)},A^{(s)})\nonumber\\
&\quad-\frac16\sum_{\substack{r+s+t=n\\ r\le s\le t}}M_{rst}\left[1-\frac{4(rs+st+tr)}{(2n+3)(n-1)}\right]
R(A^{(r)},A^{(s)},A^{(t)}),\\
\bm c_n\equiv\nabla\times\bm\Psi_n&=\sum_{\substack{r+s=n\\ r<s}}\frac{s-r}{n}\,\bm K(A^{(r)},A^{(s)}),
\end{align}
with multiplicities $M_{rs}=1$ for equal and 2 for different orders, and $M_{rst}=1$, 3 or 6 for three, two or no equal
orders. The first coefficients are $d_2=-\frac3{14}Q_{11}$, $d_3=-\frac59Q_{12}-\frac1{18}R_{111}$,
$\bm c_3=\frac13\bm K_{12}$, $d_4=-\frac7{11}Q_{13}-\frac{17}{66}Q_{22}-\frac{13}{66}R_{112}$, $\bm c_4=\frac12\bm K_{13}$.
There is no transverse part at orders 1 and 2, and from order 4 the transverse displacements of lower orders feed the
scalar sources through their full gradients. The vector is recovered by the Helmholtz inversion
$\bm\Psi_n(\vk)=[-i\vk\,d_n(\vk)+i\vk\times\bm c_n(\vk)]/k^2$; the velocity follows from
$\dd\bm\Psi/\dd\ln a=\sum_nng^n\bm\Psi_n$.

\begin{implementation}[Implementation: arbitrary order]
Products are formed with conservative full-bandwidth work meshes (widths exceed $2PK$ for order $P$ and input support
$K$), so high orders on full-support fields are very expensive: \code{plan_eds_lpt} reports the owned memory first and
\code{max_owned_bytes_per_rank} limits it. Outputs can be retained per order or streamed, and intermediate fields
cached or recomputed. The cubic mixed determinant is summed as paired $2\times2$ minors so that exactly rank-one
(plane-wave) deformations have exactly zero nonlinear sources. Empty MPI slabs are not supported; general cosmologies
and automatic SimulationCode wiring are separate work.
\begin{cppcode}
EdSLPTOptions options;
options.max_order = 6;
options.input_mode_limit = {1, 1, 1};                 // the strictly sub-Nyquist support of the seed
options.max_owned_bytes_per_rank = 512ULL << 20;
auto plan   = plan_eds_lpt(delta_ref, options, EdSLPTOutput::RetainOrders);
auto fields = compute_eds_lpt(delta_ref, options);
auto sums   = fields.evaluate(0.5, options.max_owned_bytes_per_rank);   // Psi and dPsi/dln a at g = 0.5
\end{cppcode}
\end{implementation}

\begin{validation}[Validation: LPT]
\begin{itemize}
\item The 1/2/3LPT fields (with and without padding, including the transverse third order) against independent sparse
      integer-wave-vector convolutions without wrapping (\file{tests/test_lpt.cpp}, 84 checks).
\item Arbitrary order: every Fourier coefficient against an independent Python oracle that extracts the coefficients
      of the original determinant/cofactor equations and Cauchy invariants, through order 6 for a fixture of three
      non-coplanar waves with unequal complex amplitudes (relative errors $2\times10^{-15}$ in double, $4\times10^{-7}$ in
      float); Matsubara's explicit kernels through order 4; translations, reflections, permutations, plane-wave
      exactness, and identical results serially, with OpenMP and on 1, 2 and 4 MPI ranks.
\item The example \file{LPT/example} reports per-order norms and the minimum of $\det(I+\nabla\bm\Psi)$: a negative
      Jacobian marks shell crossing, beyond which single-stream LPT is not valid.
\end{itemize}
\end{validation}

\section{Augmented LPT and spherical collapse}\label{sec:lpt-alpt}

On small scales LPT overshoots: particles stream through collapsed regions. The spherical-collapse approximation of
Bernardeau \cite{Bernardeau1994} replaces the divergence of the displacement by that of a spherical top hat,
\begin{equation}
\nabla\cdot\bm\Psi_{\rm SC}=3\left(\sqrt{\max\left(1-\tfrac23\delta(\vq,a),0\right)}-1\right),\qquad
\delta(\vq,a)=\frac{D(a)}{D(a_{\rm ini})}\delta_{\rm ini}(\vq),
\end{equation}
which equals the 1LPT divergence to linear order (\code{compute_spherical_collapse_potential}). Augmented LPT (ALPT)
\cite{KitauraHess2013} combines 2LPT on large scales with spherical collapse on small scales,
\begin{equation}
\phi_{\rm ALPT}=\phi_{\rm SC}+\mathcal S_R*\left(\phi^{(1)}+\phi^{(2)}-\phi_{\rm SC}\right),
\end{equation}
with a smoothing filter $\mathcal S_R$ (Gaussian, top-hat or sharp-$k$) of scale $R$
(\code{compute_ALPT_potential_fourier}).

\section{Particle initial conditions}\label{sec:lpt-ic}

Particle initial conditions combine the pieces of \cref{ch:growth-transfer,ch:random-fields} and this chapter:
\begin{enumerate}
\item a Gaussian (or non-Gaussian) realisation of the linear CDM+baryon density at the initial redshift, from the
      transfer provider and the growth factors (back-scaled or not);
\item the LPT potentials to order \param{ic_LPT_order} (1, 2 or 3), with each order scaled by its growth factor;
\item particles on a regular lattice (or a glass, \cref{sec:lpt-glass}) displaced by $\bm\Psi(\vq)$, interpolated from
      the grid, with velocities from \cref{eq:lpt-velocity} and the growth rates of each order.
\end{enumerate}
Two species, massive-neutrino particles and the $N$-body gauge are described in \cref{ch:growth-transfer,ch:simulation-code}.
For COLA the LPT fields are kept for the whole run. SimulationCode can also start from an existing snapshot
(\code{ic_random_field_type = read_particles} or \code{read_phases}), reconstructing the density by mass assignment
and the LPT fields from it, or exactly from the particles' displacements and velocities
(\param{ic_reconstruct_exact}).

\section{Redshift-space reconstruction}\label{sec:lpt-reconstruction}

\code{RSDReconstructionFourierMethod} removes redshift-space distortions from tracers in a periodic box with a fixed line
of sight $\hat{\bm r}$. In linear theory the displacement $\bm\Psi$ of the tracers (bias $b$, $\beta=f/b$) satisfies
\begin{equation}
\nabla\cdot(b\bm\Psi)+\beta\nabla\cdot\left[(b\bm\Psi\cdot\hat{\bm r})\hat{\bm r}\right]=-\delta_{\rm tracer},
\end{equation}
which is solved in Fourier space assuming that the second term is curl free, after smoothing the density with a filter
of scale $R$; the RSD shift $\bm\Psi_{\rm rsd}=\beta(b\bm\Psi\cdot\hat{\bm r})\hat{\bm r}$ is then subtracted. The
iterative scheme of Burden et al.\ \cite{Burden2015} starts from the observed positions $\bm s$, computes $b\bm\Psi_n$
from the density of the current estimate $\vx_{n-1}$ and sets $\vx_n=\bm s-\beta(b\bm\Psi_n(\bm s)\cdot\hat{\bm r})\hat{\bm r}$;
it converges for $\beta<1$ in 3--5 iterations. The assignment and interpolation windows are not deconvolved. Only a
periodic box with a fixed line of sight is supported here; BAO reconstruction of surveys is in \cref{ch:surveys}.
\code{particles_to_redshiftspace} applies $\bm s=\vx+c_v(\bm v\cdot\hat{\bm r})\hat{\bm r}$, with
$c_v=1/(100\,aH/H_0\,B)$ for peculiar velocities in km/s and the box $B$ in $\Mpch$ (a negative $c_v$ undoes it).

\section{Glass loads}\label{sec:lpt-glass}

A regular lattice imprints its symmetry on the early evolution; a \emph{glass}, a particle distribution with no
preferred directions and very little power on large scales, is the usual alternative. \file{GlassIC/GlassIC.h}
generates periodic equal-mass glasses by evolving a Poisson load (\code{poisson_load}, positions from Philox streams
keyed by the seed and the particle ID, so independent of the task and thread counts) under \emph{repulsive} gravity
with damping, until it settles:
\begin{itemize}
\item the force is softened periodic P$^3$M gravity with the sign reversed for both the long-range (CIC, Gaussian
      split, continuous gradient, compensation of the assignment and interpolation windows) and the short-range part
      (the spline kernel of \cref{ch:short-range}); reversing only the mesh force would leave close pairs unconstrained;
\item each step holds the acceleration fixed for an exact damped kick,
      $\bm v_{\rm next}=e^{-\gamma\Delta t}\bm v+\frac{1-e^{-\gamma\Delta t}}{\gamma}\bm a$, $\vx_{\rm next}=\vx+\Delta t\,\bm v_{\rm next}$,
      with a global velocity factor limiting the largest step to \code{max_step_spacing} times the mean spacing;
\item convergence requires the rms force and velocity below their tolerances for \code{stable_steps} consecutive steps;
      the result reports \code{Converged}, \code{MaximumSteps} or \code{NonFinite} and the history.
\end{itemize}
A converged residual alone does not make a load suitable for science: its discreteness power, anisotropy and
close-pair statistics should be checked at the intended particle number. Loads are saved in a versioned, sharded
text format (\file{LoadIO.h}) that can be read back with a different number of tasks. Perturbations are applied with
\code{apply_3lpt(load, delta_fourier, dlogDdt, dealiasing)} (EdS-factorised 3LPT sampled by CIC at the glass
positions) or \code{apply_fields} (precombined displacement and velocity fields, e.g.\ with other growth
prescriptions), and \code{density_from_provider} bridges a transfer provider to the Philox generator. Multi-species
glasses and SimulationCode integration are future work.

\begin{cppcode}[A glass load with 3LPT initial conditions]
#include <FML/GlassIC/GlassIC.h>
using namespace FML::GLASSIC;
auto load = poisson_load(32768, 12345, 250.0);   // count, seed, box [Mpc/h]
Options options;                                  // mesh = 0: smallest task multiple >= 3 N^(1/3)
options.max_steps = 3000;
auto result = relax(load, options);              // check result.status
density_from_provider(delta, provider, Species::CDMBaryon, a_ini, 250.0, field_seed, /*fixed=*/false);
apply_3lpt(load, delta, dlogDdt, FML::COSMOLOGY::LPT::Dealiasing::Padding);
\end{cppcode}
